\section{All-band conservation and the reassembly barrier}
\label{sec:tpc48-conservation}

The output-side square function is exact and useful bookkeeping.  It
must not be confused with a new atomic estimate.

\subsection{Amplitude and diagonal conservation}

Let \(Sx=\sum_{e\in\cE}x_eu_e\) be any finite synthesis in
\(\widetilde\Hh\), and define
\begin{equation}
 \cD(x):=\sum_{e\in\cE}|x_e|^2\|u_e\|^2.
 \label{eq:tpc48-abstract-diagonal}
\end{equation}
For each orbit band, put
\begin{equation}
 S_kx:=\Pi_kSx,
 \qquad
 \cD_k(x):=\sum_e|x_e|^2\|\Pi_ku_e\|^2.
 \label{eq:tpc48-band-diagonal}
\end{equation}

\begin{proposition}[Exact all-band conservation]
\label{prop:tpc48-all-band-conservation}
For every \(x\),
\begin{equation}
 \sum_k\|S_kx\|^2=\|Sx\|^2,
 \qquad
 \sum_k\cD_k(x)=\cD(x).
 \label{eq:tpc48-all-band-conservation}
\end{equation}
Consequently, if
\(\|S_kx\|^2\le C\cD_k(x)\) for every \(k\) with one common constant
       \(C\), then \(\|Sx\|^2\le C\cD(x)\).  Conversely, whenever
\(\cD(x)>0\), at least one band with \(\cD_k(x)>0\) satisfies
\begin{equation}
 \frac{\|S_kx\|^2}{\cD_k(x)}
 \ge\frac{\|Sx\|^2}{\cD(x)}.
 \label{eq:tpc48-one-band-inherits-ratio}
\end{equation}
\end{proposition}

\begin{proof}
The first identity is Parseval for the orthogonal resolution
\eqref{eq:tpc48-projector-resolution}.  Apply the same identity to each
atom \(u_e\), multiply by \(|x_e|^2\), and sum to obtain the diagonal
identity.  Summing the bandwise inequalities gives the first
consequence.  The full ratio is the \(\cD_k\)-weighted average of the
nonzero band ratios, which proves the second.
\end{proof}

Thus an output tiling cannot hide a pre-existing coherent-box or
tensor-to-atomic obstruction: some band inherits at least the full
ratio.  The conclusion is abstract and does not say that every actual
band is bad.

\subsection{Collapse of the difference channels}

Let
\begin{equation}
 \cQ_L:\bigoplus_{q\bmod L}\widetilde\Hh
 \longrightarrow\widetilde\Hh,
 \qquad
 \cQ_L((z_q)_q):=\sum_qz_q.
 \label{eq:tpc48-channel-collapse}
\end{equation}
Then \(\|\cQ_L\|^2=L\).  The channel lift factors the unlabelled
determinant synthesis:
\begin{equation}
 \cQ_L\cA_{\alpha,L}\bm y
 =\sum_{m\in\cM}U_{m\alpha}y_m
 =:T_\alpha\bm y.
 \label{eq:tpc48-channel-factorization}
\end{equation}

\begin{theorem}[Sharp unrestricted reassembly barrier]
\label{thm:tpc48-unrestricted-barrier}
For every \(\alpha\in\T\),
\begin{equation}
 \boxed{\|T_\alpha\|^2=K.}
 \label{eq:tpc48-unrestricted-barrier}
\end{equation}
Hence
\begin{equation}
 K\le L R_L(\alpha;\cM),
 \label{eq:tpc48-factor-product}
\end{equation}
and no universal inequality for arbitrary independent orbit inputs can
replace the factor \(K\) by a smaller all-band constant.
\end{theorem}

\begin{proof}
Cauchy--Schwarz gives
\(\|\sum_mU_{m\alpha}y_m\|^2\le K\sum_m\|y_m\|^2\).
For a nonzero finitely supported \(g\), take
\(y_m=U_{m\alpha}^{-1}g\).  Then the output is \(Kg\), while the input
norm squared is \(K\|g\|^2\), proving equality.  The product bound is
also immediate from
\(T_\alpha=\cQ_L\cA_{\alpha,L}\) and
\cref{thm:tpc48-channel-norm}.
\end{proof}

The witness is allowed in the ambient arbitrary-input class and can be
finitely supported in the orbit coordinate.  It is not asserted to be
an actual source-mask coefficient.  Its role is to identify exactly
which extra structure a positive continuation must use.

\subsection{A no-wrap collision bound}

On the shift range relevant below, the slope points do not wrap around
the quotient circle.

\begin{proposition}[Elementary tube-load bounds]
\label{prop:tpc48-load-bounds}
Assume \(\cM\subset[M_-,M_+]\cap\Z\), \(0<t<1\),
\(L\le M_+\), and
\begin{equation}
 H_0M_+t<1.
 \label{eq:tpc48-load-no-wrap}
\end{equation}
Then
\begin{equation}
 \left\lceil\frac KL\right\rceil
 \le R_L(t\pmod1;\cM)
 \le\min\left\{K,1+\frac2{LH_0t}\right\}.
 \label{eq:tpc48-load-bounds}
\end{equation}
If, more explicitly,
\begin{equation}
 \cM\subset[M,2M]\cap\Z,
 \qquad K\asymp M,
 \qquad L\le M,
 \label{eq:tpc48-dense-slope-box}
\end{equation}
and
\begin{equation}
 \frac aM\le t\le\frac bM,
 \qquad 0<a<b<\frac1{2H_0},
 \label{eq:tpc48-coarse-edge-slab}
\end{equation}
then
\begin{equation}
 R_L(t\pmod1;\cM)\asymp_{a,b,H_0}\frac ML.
 \label{eq:tpc48-dense-load}
\end{equation}
\end{proposition}

\begin{proof}
The lower bound is \eqref{eq:tpc48-load-pigeonhole}.  Under
\eqref{eq:tpc48-load-no-wrap}, the points \(H_0mt\) are separated by
at least \(H_0t\) on an interval of length less than one.  A quotient
cell has length \(1/L\), so it contains at most
\(1+2(LH_0t)^{-1}\) such points; the factor two harmlessly covers a
cell meeting the chosen cut of the quotient circle.  This is the upper
bound.  On
\eqref{eq:tpc48-coarse-edge-slab}, the upper bound is \(O(M/L)\), while
the dense-slope hypothesis and the pigeonhole lower bound give the
matching estimate.
\end{proof}

For the critical choice \(L\asymp J\), a dense slope box therefore has
channel load \(M/J\) on the coarse-edge slab.  However, multiplying by
the \(L\)-channel collapse returns \(M\), exactly as
\eqref{eq:tpc48-factor-product} predicts.
